\chapter{Розвідка загроз, індикатори компрометації та атрибуція}
\chaptershort{Розвідка загроз та атрибуція}
\label{chap:ch12}

Виявлення й аналіз окремих зразків (розділи~\ref{chap:ch06}--\ref{chap:ch11}) дають фрагменти, які \term{розвідка загроз} (threat intelligence) об'єднує в цілісне, контекстуалізоване знання про супротивника: хто атакує, якими засобами, з якою метою, за якими ознаками їх розпізнати й передбачити. Розвідка загроз перетворює реактивне виявлення на проактивний захист: знаючи тактики, техніки й інфраструктуру ймовірних супротивників, захисник готується заздалегідь. Цей розділ розглядає поняття й рівні розвідки загроз, індикатори компрометації та їх життєвий цикл, стандарти обміну, моделі аналізу й атрибуції, її можливості й межі, а також роль CERT-UA в Україні, і завершується кейсом на основі кампаній WannaCry та NotPetya. Розвідка загроз живить інженерію виявлення (розділ~\ref{chap:ch11}), полювання за загрозами (розділ~\ref{chap:ch13}), реагування (розділ~\ref{chap:ch16}) і стратегічне управління (розділ~\ref{chap:ch17}).

\section{Поняття й рівні розвідки загроз}\label{sec:ch12-levels}

\term{Розвідка загроз} --- це засноване на доказах знання про наявні чи ймовірні загрози, придатне для прийняття рішень~\cite{ch12:mcmillan}. Ключові слова --- \emph{доказовість} (спирається на спостереження, а не припущення) і \emph{придатність для рішень} (не дані заради даних, а знання, що змінює дії). Розрізняють три рівні за горизонтом і аудиторією (табл.~\thechapter.1): \emph{тактичний} (технічні індикатори й техніки для негайного виявлення --- для аналітиків і засобів захисту); \emph{оперативний} (кампанії, інструменти, інфраструктура конкретних угруповань --- для реагування й полювання); \emph{стратегічний} (мотиви, можливості, тенденції супротивників --- для керівництва й стратегії, розділ~\ref{chap:ch17}).

\tabhead{Таблиця~\thechapter.1 --- Рівні розвідки загроз}

\begin{center}\footnotesize
\begin{tabularx}{\linewidth}{L{20mm} Y L{20mm}}
\toprule
Рівень & Зміст & Аудиторія \\
\midrule
Тактичний & індикатори, техніки & аналітики, засоби \\
Оперативний & кампанії, інфраструктура & реагування, полювання \\
Стратегічний & мотиви, тенденції & керівництво \\
\bottomrule
\end{tabularx}
\end{center}

Цінність розвідки визначається не обсягом, а \emph{релевантністю} й \emph{своєчасністю}: індикатор, актуальний для моделі загроз організації (розділ~\ref{chap:ch04}) і отриманий вчасно, вартий більше за гори нерелевантних даних. Тому зріла розвідка починається з визначення пріоритетних розвідувальних вимог --- які саме загрози важливі для цієї організації.

Розвідка загроз проходить класичний \emph{розвідувальний цикл}: планування (визначення вимог) $\to$ збір (із власної телеметрії, обміну, відкритих і закритих джерел) $\to$ оброблення (нормалізація, дедублікація, збагачення) $\to$ аналіз (перетворення даних на знання) $\to$ поширення (доведення до споживачів у придатній формі) $\to$ зворотний зв'язок (оцінка корисності й уточнення вимог). Пропуск будь-якого етапу знижує цінність: збір без планування дає нерелевантні гори даних; аналіз без поширення лишає знання невикористаним; поширення без зворотного зв'язку не дозволяє вдосконалюватися. Саме дисциплінований цикл, а не доступ до «стрічок індикаторів», відрізняє зрілу розвідкову функцію від простого споживання чужих даних.

Важливо розрізняти \emph{споживання} та \emph{виробництво} розвідки. Більшість організацій починають зі споживання: підписуються на зовнішні канали індикаторів і звітів. Зрілі організації додають виробництво: аналізують власні інциденти (розділ~\ref{chap:ch16}), видобувають із них індикатори й TTP, збагачують і повертають у спільноту. Виробництво цінне тим, що власна телеметрія відображає саме ті загрози, які атакують цю організацію, --- найрелевантніше джерело. Поєднання споживання й виробництва замикає контур: зовнішня розвідка спрямовує увагу, а внутрішній аналіз уточнює й доповнює її під власний контекст.

\section{Індикатори компрометації та їх життєвий цикл}\label{sec:ch12-ioc}

\term{Індикатор компрометації} (IoC) --- спостережувана ознака, наявність якої свідчить про ймовірну компрометацію: геш шкідливого файлу, адреса чи домен командного сервера, ключ реєстру персистентності, патерн мережевого трафіку, рядок у пам'яті. IoC --- робочий матеріал тактичної розвідки: ними наповнюють правила виявлення (розділ~\ref{chap:ch11}), їх перевіряють у телеметрії (розділ~\ref{chap:ch10}).

\begin{tikzfig}
\begin{tikzpicture}[font=\scriptsize,
  b/.style={draw,rounded corners=2pt,minimum height=5.5mm,text width=17mm,align=center,font=\tiny},
  >={Stealth[length=1.4mm]}]
\node[b] (c) at (0,0) {збір (аналіз, обмін)};
\node[b] (v) at (2.2,0) {валідація};
\node[b] (d) at (4.4,0) {розгортання (виявлення)};
\node[b] (e) at (6.6,0) {збагачення};
\node[b] (r) at (4.4,-1.3) {старіння / вивід};
\draw[->] (c)--(v); \draw[->] (v)--(d); \draw[->] (d)--(e);
\draw[->] (e) to[out=-90,in=0] (r); \draw[->] (r)--(d.south);
\end{tikzpicture}
\figcap{Рис.~\thechapter.1 --- Життєвий цикл індикатора компрометації}
\end{tikzfig}

IoC мають \term{життєвий цикл}: збір (з власного аналізу, обміну, відкритих джерел) $\to$ валідація (перевірка достовірності й релевантності) $\to$ розгортання $\to$ збагачення контекстом $\to$ старіння й виведення. Критична властивість --- \emph{старіння}: адреси й домени зловмисник швидко змінює (нижні рівні піраміди болю, розділ~\ref{chap:ch11}), тому старий IoC дає хибні тривоги (адресу перепризначено легітимному сервісу) або пропуски. Тому IoC без терміну придатності й контексту --- радше шум, ніж розвідка. Це знову підкреслює перевагу стійких поведінкових індикаторів (TTP) над ефемерними атомарними.

Старіння індикаторів зручно моделювати спаданням їх розвідувальної цінності в часі. Якщо початкову цінність свіжого індикатора позначити $v_0$, а темп старіння --- сталою $\lambda$, характерною для типу індикатора, то цінність через час $t$ після появи

\begin{equation}
v(t) = v_0\,e^{-\lambda t},
\label{eq:ch12-decay}
\end{equation}

де велике $\lambda$ (швидке старіння) відповідає ефемерним індикаторам нижніх рівнів піраміди болю (адресам, доменам), а мале $\lambda$ --- стійким поведінковим ознакам (TTP). Формула~\eqref{eq:ch12-decay} формалізує інтуїцію про «термін придатності» й обґрунтовує диференційовану політику: атомарні індикатори потребують частого оновлення й швидкого виведення, тоді як поведінкові зберігають цінність надовго. Половинний час життя індикатора $t_{1/2} = \ln 2 / \lambda$ слугує практичним орієнтиром для планування перегляду бази індикаторів.

\section{Стандарти обміну розвідкою}\label{sec:ch12-sharing}

Цінність розвідки зростає з обміном: загроза, виявлена в одній організації, застерігає решту. Для машиночитаного обміну розроблено стандарти: \term{STIX}~\cite{stix} --- структурований формат опису загроз (індикатори, техніки, угруповання, кампанії та зв'язки між ними); \term{TAXII}~\cite{ch12:taxii} --- протокол їх передавання. Платформи обміну, зокрема відкрита MISP~\cite{ch12:misp}, автоматизують накопичення, кореляцію й поширення розвідки між спільнотами й галузевими центрами.

Обмін має бар'єри: довіра (чи достовірне джерело), конфіденційність (індикатор може розкрити факт компрометації його джерела), швидкість (застарілий індикатор марний), якість (без контексту й валідації обмін поширює шум). Для керування чутливістю застосовують протоколи маркування (на кшталт TLP --- traffic light protocol), що визначають, кому можна передавати інформацію. Попри бар'єри, колективний обмін є силовим множником: він зміщує асиметрію на користь захисників, дозволяючи вчитися на чужому досвіді, перш ніж зазнати тієї самої атаки.

Структуроване машиночитане подання розвідки має й самостійну цінність поза передаванням. Коли індикатори, техніки, інструменти, угруповання й кампанії описані в єдиній формальній моделі (як у STIX) зі зв'язками між сутностями, розвідка стає \emph{графом знань}, що піддається автоматизованому аналізу: можна запитати, які кампанії застосовували даний інструмент, яку інфраструктуру поділяють два угруповання, які жертви зазнали схожих TTP. Це перетворює розрізнені звіти на пов'язану базу, у якій нове спостереження автоматично зіставляється з накопиченим і або підтверджує відому кампанію, або сигналізує про нову. Без такої структуризації розвідка лишається купою текстів, корисність яких обмежена здатністю окремого аналітика все це пам'ятати.

\section{Моделі аналізу й атрибуція}\label{sec:ch12-attribution}

Для структурування аналізу застосовують моделі. \term{Діамантова модель}~\cite{ch12:diamond} описує подію через чотири вершини --- супротивник, можливості (інструменти), інфраструктура, жертва --- та зв'язки між ними, що допомагає від відомих елементів виводити невідомі (наприклад, від інфраструктури --- до інших жертв). У поєднанні з ланцюгом ураження й ATT\&CK (розділ~\ref{chap:ch04}) вона дає каркас для накопичення знань про угруповання.

\begin{tikzfig}
\begin{tikzpicture}[font=\scriptsize,
  n/.style={draw,rounded corners=2pt,minimum height=5mm,text width=20mm,align=center,font=\tiny},
  >={Stealth[length=1.3mm]}]
\node[n] (adv) at (0,1.3) {супротивник};
\node[n] (vic) at (0,-1.3) {жертва};
\node[n] (cap) at (-2.3,0) {можливості (інструменти)};
\node[n] (inf) at (2.3,0) {інфраструктура};
\draw[-] (adv)--(cap); \draw[-] (adv)--(inf);
\draw[-] (vic)--(cap); \draw[-] (vic)--(inf);
\draw[-] (adv)--(vic); \draw[-] (cap)--(inf);
\end{tikzpicture}
\figcap{Рис.~\thechapter.2 --- Діамантова модель вторгнення: чотири вершини та зв'язки, що дозволяють виводити невідомі елементи від відомих}
\end{tikzfig}

\term{Атрибуція} --- віднесення атаки до конкретного актора --- є одним із найскладніших і найделікатніших завдань~\cite{ch12:rid}. Вона спирається на сукупність свідчень: технічні (інфраструктура, повторно використаний код, артефакти збірки, мовні сліди, робочий час за часовим поясом), оперативні (вибір цілей, TTP) і контекстуальні (геополітичний мотив). Жодне свідчення саме не є доказом, тому атрибуція --- це оцінка ймовірності за сукупністю, а не однозначне твердження. Кластеризацію активності в угруповання позначають умовними іменами (APT-номери, кодові назви) ще до встановлення реального актора.

Доказовий характер атрибуції зручно виразити через накопичення свідчень. Якщо кожне незалежне свідчення $i$ підтримує гіпотезу про актора з вагою (інформативністю) $w_i$ та спостерігається (значення індикатора $x_i\in\{0,1\}$), то сукупну впевненість можна подати нормованою зваженою сумою

\begin{equation}
C = \frac{\sum_i w_i x_i}{\sum_i w_i} \in [0;1],
\label{eq:ch12-confidence}
\end{equation}

що зростає з кількістю й вагомістю узгоджених свідчень. Формула~\eqref{eq:ch12-confidence} підкреслює дві речі: жодне окреме свідчення ($x_i=1$ за малого $w_i$) не дає впевненості, а висока впевненість потребує сукупності незалежних свідчень; крім того, вагу $w_i$ слід занижувати для ознак, які легко підробити (хибні прапори), --- саме тому мовні сліди чи чужі інструменти важать менше за складно підроблювані оперативні патерни.

Межі атрибуції принципові. По-перше, \emph{хибні прапори}~\cite{ch12:attribution}: зловмисник навмисно залишає сліди, що вказують на іншого (чужу мову, чужі інструменти, чужу інфраструктуру), аби ввести в оману. По-друге, повторне використання інструментів і витоки коду розмивають зв'язок «інструмент~$\to$~актор». По-третє, атрибуція має політичні наслідки, тому потребує особливо високого стандарту доказовості. Для захисту точна атрибуція до держави чи особи часто й не потрібна: достатньо стійко розпізнавати \emph{угруповання} за його TTP, щоб виявляти й випереджати його дії, незалежно від того, хто за ним стоїть.

\section{Розвідка загроз в Україні: CERT-UA}\label{sec:ch12-ukraine}

Україна --- полігон і ціль найактивніших у світі кампаній ЗПК, тому національна розвідка загроз має критичне значення. \term{CERT-UA}~\cite{cert-ua} --- урядова команда реагування на комп'ютерні надзвичайні події у складі Держспецзв'язку (розділ~\ref{chap:ch02}) --- збирає, аналізує й поширює індикатори та аналітичні звіти про атаки на українські організації, координує реагування й міжнародний обмін. Через безпрецедентну інтенсивність атак на критичну інфраструктуру (розділ~\ref{chap:ch14}) український досвід і розвідка загроз мають світове значення: багато кампаній, уперше виявлених в Україні, згодом загрожували всьому світу. Вітчизняні дослідження систематизують застосування розвідки загроз для проактивного захисту~\cite{ch12:ti-ua} та діяльність APT-угруповань проти критичної інфраструктури України~\cite{ch12:apt-ua}.

Український контекст вирізняється кількома особливостями, важливими для розуміння розвідки загроз. По-перше, тут спостерігається рідкісна концентрація діяльності найспроможніших державних угруповань, що застосовують повний спектр --- від шпигунства до деструктивних атак на енергетику (розділ~\ref{chap:ch14}). По-друге, саме в Україні вперше спостерігали техніки, які згодом поширилися світом (деструктивні вайпери під виглядом вимагачів, атаки на промислову автоматику), що робить українську розвідку випереджальним індикатором для глобальної спільноти. По-третє, висока інтенсивність атак прискорила розвиток практик швидкого обміну індикаторами між державним сектором, критичною інфраструктурою та міжнародними партнерами. Для аспіранта-дослідника це означає, що українські дані --- не локальний окремий випадок, а один із найбагатших у світі емпіричних матеріалів для вивчення сучасних загроз і методів протидії.

\section{Кейс: від зразка до розвідки (WannaCry, NotPetya)}\label{sec:ch12-case}

Розглянемо, як окремі спостереження перетворюються на розвідку, на прикладі двох кампаній 2017~року, що зачепили Україну й світ (захисна/аналітична задача).

\emph{Крок 1. Технічний аналіз.} Статичний і динамічний аналіз (розділи~\ref{chap:ch06}--\ref{chap:ch07}) зразків WannaCry~\cite{ch12:wannacry} і NotPetya~\cite{ch01:greenberg} дає технічні деталі: механізм поширення (експлуатація мережевої вразливості), шифрування, індикатори. Виявляється, що NotPetya, попри вигляд вимагача, був деструктивним вайпером без реальної можливості відновлення --- висновок, можливий лише через аналіз, а не з самозаявленої поведінки.

\emph{Крок 2. Індикатори.} Формуються IoC: геші, мережеві патерни, домени, ознаки поведінки --- для негайного виявлення (розділ~\ref{chap:ch11}).

\emph{Крок 3. Кампанія й контекст.} Діамантова модель~\cite{ch12:diamond} і ATT\&CK зводять спостереження в картину кампанії: вектор початкового доступу (для NotPetya --- компрометація оновлення бухгалтерського ПЗ, тобто ланцюг постачання, розділ~\ref{chap:ch15}), цілі, географія, TTP.

\emph{Крок 4. Атрибуція.} За сукупністю технічних і контекстуальних свідчень активність кластеризують і відносять до угруповання, з урахуванням ризику хибних прапорів~\cite{ch12:attribution}. Для NotPetya вибір цілей (переважно українські організації) і контекст є вагомими оперативними свідченнями.

\emph{Крок 5. Поширення й застосування.} Розвідку (STIX/TAXII, через CERT-UA й спільноти) поширюють, щоб решта організацій застосувала її проактивно: оновила вразливе ПЗ, розгорнула виявлення, перевірила телеметрію ретроспективно.

Кейс показує повний шлях від зразка до дієвої розвідки й ілюструє ключові уроки: самозаявлена поведінка (вимагач) може маскувати справжню мету (деструкція); вектор ланцюга постачання робить навіть оновлення недовіреним; а своєчасний обмін розвідкою міг би обмежити глобальне поширення.

\section*{Висновки до розділу}
\addcontentsline{toc}{section}{Висновки до розділу}

Розвідка загроз об'єднує фрагменти виявлення й аналізу в контекстуалізоване знання про супротивника, перетворюючи реактивний захист на проактивний. Вона має три рівні --- тактичний, оперативний, стратегічний (табл.~\thechapter.1), --- а її цінність визначають доказовість, релевантність і своєчасність, а не обсяг. Індикатори компрометації --- робочий матеріал тактичної розвідки --- мають життєвий цикл зі старінням, через що ефемерні атомарні IoC поступаються стійким поведінковим TTP. Обмін розвідкою (STIX, TAXII, MISP) є силовим множником, що зміщує асиметрію на користь захисників, попри бар'єри довіри, конфіденційності та якості. Моделі аналізу (діамантова, ATT\&CK) структурують накопичення знань про угруповання; атрибуція спирається на сукупність свідчень як оцінка ймовірності й обмежена хибними прапорами та повторним використанням інструментів, тому для захисту часто достатньо розпізнавати угруповання за TTP. В Україні CERT-UA виконує критичну роль національної розвідки загроз, а український досвід має світове значення. Кейс WannaCry/NotPetya демонструє повний шлях від зразка до дієвої розвідки й уроки маскування мети, вектора ланцюга постачання та цінності своєчасного обміну. Головний висновок: розвідка загроз --- це міст від технічного аналізу до стратегічних рішень, а її сила --- у релевантності, контексті, стійких індикаторах і колективному обміні.

\selfcheck
\begin{enumerate}
\item Що таке розвідка загроз? Поясніть ключові ознаки доказовості й придатності для рішень.
\item Опишіть три рівні розвідки загроз (табл.~\thechapter.1) за горизонтом і аудиторією.
\item Чому цінність розвідки визначається релевантністю й своєчасністю, а не обсягом?
\item Що таке індикатор компрометації? Наведіть приклади й опишіть їх життєвий цикл.
\item Чому старіння IoC є критичною проблемою? Як це пов'язано з пірамідою болю (розділ~\ref{chap:ch11})?
\item Які стандарти обміну розвідкою існують (STIX, TAXII, MISP)? Яку задачу розв'язує кожен?
\item Назвіть бар'єри обміну розвідкою. Чому попри них обмін є силовим множником?
\item Опишіть діамантову модель аналізу вторгнень. Як вона допомагає виводити невідомі елементи?
\item Що таке атрибуція? На які групи свідчень вона спирається й чому є оцінкою ймовірності?
\item Які принципові межі має атрибуція? Що таке хибні прапори? Чому для захисту достатньо розпізнавання угруповання за TTP?
\item Яку роль відіграє CERT-UA? Чому український досвід розвідки загроз має світове значення?
\item Опишіть шлях від зразка до розвідки на прикладі WannaCry/NotPetya (кейс підрозділу~\ref{sec:ch12-case}). Чому самозаявлена поведінка може маскувати мету?
\item Як розвідка загроз пов'язана з інженерією виявлення, реагуванням і стратегічним управлінням?
\item Дослідницьке: побудуйте профіль обраного APT-угруповання за відкритими джерелами, застосувавши діамантову модель і ATT\&CK, та оцініть повноту й суперечливість атрибуції.
\item Дослідницьке: оцініть швидкість старіння різних типів IoC (гешів, доменів, TTP) на реальному потоці розвідки й обґрунтуйте політику їх виведення з експлуатації.
\end{enumerate}
